% General CV, simple monochrome layout requested by the user.
% The existing filename is retained for the open document editor.
% Candidate facts: local evidence profile, source documents and user clarifications.
\documentclass[11pt,a4paper]{article}
\usepackage{fontspec}
\setmainfont{Times New Roman}
\usepackage[margin=0.7in]{geometry}
\usepackage{enumitem}
\usepackage{titlesec}
\usepackage{etoolbox}
\usepackage{needspace}
\usepackage[hidelinks]{hyperref}
\hypersetup{
  pdftitle={Ardhya Novri Ramadhan - CV},pdfauthor={Ardhya Novri Ramadhan},
  pdfpagemode=UseNone}
\pagestyle{plain}
\setlength{\parindent}{0pt}
\setlength{\parskip}{2pt}
\setlist[itemize]{leftmargin=13pt,itemsep=1pt,topsep=2pt,parsep=0pt}
\titleformat{\section}{\normalfont\bfseries\fontsize{11}{13}\selectfont}{}{0pt}{\MakeUppercase}[\titlerule]
\titlespacing*{\section}{0pt}{9pt}{4pt}
\newcommand{\cventry}[6]{%
  \Needspace{4\baselineskip}%
  \textbf{#3}\hfill #1\par
  \textit{#2}\ifstrempty{#4}{}{\hfill #4}\par
  #6\par\vspace{4pt}}
\newcommand{\cvitem}[2]{\textbf{#1}: #2\par}
\begin{document}
\begin{center}
{\fontsize{19}{22}\selectfont\bfseries Ardhya Novri Ramadhan}\par
\vspace{3pt}
{\small Bekasi, Indonesia\enspace |\enspace +62 812 9155 2970\enspace |\enspace
\href{mailto:ardhyasan@gmail.com}{ardhyasan@gmail.com}}\par
{\small \href{https://www.linkedin.com/in/dhya/}{linkedin.com/in/dhya}\enspace |\enspace
\href{https://ardhyanovri.my.id}{ardhyanovri.my.id}\enspace |\enspace
\href{https://github.com/vanovri}{github.com/vanovri}}
\end{center}
\raggedright

Bachelor of Engineering in Electrical Engineering from Diponegoro University, with experience
in electronics, airport electrical systems, student organizations and community projects. Develops embedded monitoring systems and
AI-assisted software, and is the sole developer of ENVIID's company website.

\section{Education}
\cventry{Aug 2022 - Sep 2026}{Bachelor of Engineering in Electrical Engineering}{Diponegoro University}{}{}{
GPA: 3.43/4.00. Graduated on 29 September 2026.}
\cventry{Jul 2019 - May 2022}{Natural Sciences}{SMA Negeri 5 Bekasi}{}{}{}

\section{Professional and Community Experience}
% The user confirmed the role started when the website was made, 12 October 2026.
\cventry{Oct 2026 - Present}{Chief Technology Officer / IT \& Digital Lead}{ENVIID Carbon Solutions}{}{}{
Sole developer of the company website, built with AI assistance, including project information,
a carbon calculator and a digital MRV workflow demonstration.}
\cventry{Jan 2026 - Feb 2026}{Community Service Participant (KKN)}{Diponegoro University}{Mutihan, Klaten}{}{
\begin{itemize}
\item Helped assemble and install a solar-powered Smart Owl House, including sensors, battery and controller. Prepared operating and maintenance guidance.
\item Assisted village administration, Posyandu activities, BUMDes bookkeeping and community cleanups. Joined mosque and cemetery painting activities.
\item Supported school scouting activities and earthquake preparedness education. Helped produce and hand over the village profile video.
\end{itemize}}
\cventry{Jul 2025 - Aug 2025}{Internship Trainee}{PT Angkasa Pura Indonesia}{}{}{
\begin{itemize}
\item Rotated for five weeks in UPS/converter systems and three weeks in electrical protection at Soekarno-Hatta International Airport.
\item Visited substations and focused on learning protection relays. Received SCADA instruction and observed PLC installations.
\end{itemize}}
\cventry{2025}{Basic Electronics Practicum Assistant}{Diponegoro University}{}{}{
Assisted students with laboratory activities and basic electronic components.}
\cventry{Jan 2025 - Feb 2025}{Electronics Technician}{Redesma Technologies}{}{}{
Worked on IoT temperature monitoring and DC motor control. Co-authored an SOP
for comparing Blynk and Thinger.io using ESP32 and a DS18B20 sensor.}

\section{Skills}
\cvitem{Electrical and embedded}{Electronics, sensor integration, C, STM32 HAL,
STM32CubeIDE, ADC, DMA, timers and UART. Field learning in UPS/converters, protection relays and SCADA.}
\cvitem{Software and media}{Python, Tkinter, React, FastAPI, Git, automated tests,
digital filtering, RMS/MAV extraction, CSV/JSON processing and Canva.}

\newpage
\section{Organizations and Event Committees}
\cventry{2025}{Head of Media Division}{Society of Electronics Engineering Division (SEED)}{}{}{
Served as Head of Media for the 2025 student bureau at Diponegoro University.}
\cventry{Class of 2025}{Google Student Ambassador}{Google Student Ambassador Program}{}{}{
Completed the student program and received recognition for championing Google AI within the campus community.}
\cventry{Sep 2024 - Dec 2025}{Student Member}{Society of Petroleum Engineers (SPE)}{}{}{
Joined professional learning activities, including technical writing and geothermal drilling sessions.}
\cventry{Oct 2024}{Workshop Instructor}{ENIGMA Vol. 3, SEED / HME UNDIP}{}{}{
Taught basic electronics workshop activities involving circuits, wiring and PCB work.}
\cventry{Sep 2024}{Organizational Delegate}{Indonesia International Sustainability Forum (ISF)}{}{}{
Participated in the sustainability forum as an organizational delegate.}
\cventry{2024}{Creative Design Staff, Kantor Media Informasi (KMI)}{Himpunan Mahasiswa Elektro, Diponegoro University}{}{}{
Supported student organization media and design work. Assigned as a co-responsible staff member for the HME Rewind documentation program.}
\cventry{2023}{Event Committee Member}{LKMM-PD, HME UNDIP}{}{}{
Served on the committee for introductory student management and leadership training.}
% Two official committee lists establish 2023; an older self-reported form says 2022.
\cventry{2023}{Logistics / Equipment Committee Member}{PMB Teknik Elektro, Diponegoro University}{}{}{
Prepared equipment and facilitated the needs of the new-student welcoming activities.}
\cventry{Nov 2022 - Jan 2023}{Campaign Contributor}{Decarbonizing Indonesia}{}{}{
Created social media campaign and design materials about decarbonization.}
\cventry{Aug 2022}{Participant}{PKKMB and Pendikar, Diponegoro University}{}{}{
Completed the 40-hour new-student orientation and character education program, 15-20 August 2022.}
\cventry{Jan 2021 - Dec 2021}{Vice Secretary}{Information and Communication Technology (ICT), SMA Negeri 5 Bekasi}{}{}{
Prepared work-program and accountability documents, and maintained attendance and member assessment records.}
\cventry{Nov 2021}{Committee Secretary and Session Instructor}{ICT Leadership Regeneration, SMA Negeri 5 Bekasi}{}{}{
Prepared event administration and delivered the secretarial training session for incoming ICT committee members.}
\cventry{2021}{Event Secretary}{CINOTIC, SMA Negeri 5 Bekasi}{}{}{
Supported proposals, reporting and event administration for the Science Techno \& Journalistic Competition.}

\newpage
\section{Projects and Technical Proposals}
% Each project keeps one dated heading and its description, avoiding redundant role lines.
\renewcommand{\cventry}[6]{%
  \Needspace{3\baselineskip}%
  \textbf{#3}\hfill #1\par
  #6\par\vspace{5pt}}
% Software implementation used AI coding assistance. Proposal outcomes are not field deployments.
\cventry{Oct 2026}{Code-repair Verification Tool}{\href{https://github.com/vanovri/repairgate}{RepairGate}}{}{}{
Developed a code-repair verifier with protected tests, independent holdout checks and exported patch decisions.
Includes an optional local Ollama adapter and a recorded-candidate demonstration.}
\cventry{Oct 2026}{Evidence-based Application Workspace}{\href{https://github.com/vanovri/applyproof}{ApplyProof}}{}{}{
Developed an application workspace with a reviewed fact ledger, explicit skill comparison, office/worksite distinction
and traceable application drafts. Optional AI selects fact IDs for verbatim drafting.}
\cventry{Oct 2026}{Structured Data Repair Tool}{\href{https://github.com/vanovri/datarescue}{DataRescue}}{}{}{
Built a local CSV/JSON validation workbench and CLI with rule-based normalization, quarantine of rejected rows and audit-report exports, covered by automated tests.}
\cventry{Oct 2026}{Serial Parser Test Toolkit}{\href{https://github.com/vanovri/uart-fixture-kit}{uart-fixture-kit}}{}{}{
Developed reusable synthetic UART fixtures and virtual-time replay to reproduce fragmented,
malformed and concatenated byte streams without serial hardware.}
\cventry{2026}{Engineering Portfolio and Browser Game}{\href{https://ardhyanovri.my.id}{Personal Portfolio / Anak Kos Survival}}{}{}{
Developed an engineering portfolio with project evidence and a simulated EMG preview.
Earlier work includes a seven-day browser game and an interactive portfolio world, retained in the project archive.}
\cventry{2026}{Solar UAV Concept and Interactive Demonstration}{Nusantara SolarWing}{}{}{
Prepared a model-based concept package using historical solar data, wind sensitivity
and battery reserve planning for maritime patrol scenarios.}
\cventry{Thesis, 2026}{Academic Thesis}{EMG Acquisition and Monitoring System}{}{}{
Academic thesis, Diponegoro University.
\begin{itemize}
\item Implemented 12-bit ADC acquisition with DMA on STM32F411CEU6 and SEN0240 at a nominal 4 kHz sampling rate.
\item Applied 20 Hz high-pass, 500 Hz low-pass and 50 Hz notch filtering, with RMS/MAV extraction over 400 samples.
\item Built the Python/Tkinter interface for UART telemetry at 115200 baud, live plots, status monitoring and CSV export, with coding-agent assistance. Teammates handled the stimulator and fuzzy control.
\end{itemize}}
\cventry{2025}{Engineering Economics Group Assignment}{PLTA Jatigede Case Study}{}{}{
Contributed the introduction and task allocation to a group presentation covering
NPV, benefit-cost ratio and payback period.}
\cventry{2025}{Technology Development Proposal}{Redesma Technologies / Timor-Leste Concept}{}{}{
Authored a proposal for rural technology options, including solar-powered irrigation,
community warning systems and digital monitoring.}
\cventry{2024}{Environmental Technology Concept}{EnvironmentID: IoT-powered Biofiltration}{}{}{
Co-authored an abstract proposing smart sensors and oyster ecosystems for coastal water-quality monitoring and biofiltration.}
\cventry{2024}{Public Service Technology Proposal}{GOVTECH / Gov-AI Hackathon Proposal}{}{}{
Prepared a proposal for integrated digital public services using APIs, responsive interfaces and AI-assisted information services.}

\newpage
\section{Courses, Workshops and Certifications}
\cvitem{Sep 2026}{\textbf{Introduction to Internet of Things}, Cisco Networking Academy. Completed 27 September 2026.}
\cvitem{Sep 2026}{\textbf{AI Fundamentals}, Google via Coursera. Completed 26 September 2026.}
\cvitem{Sep 2025}{\textbf{SPENTRODUCE: Exploring the Downstream Oil and Gas Industry}, SPE UNDIP Student Chapter. Participant.}
\cvitem{Jan 2025}{\textbf{AI-Powered Shopping Ads Certification}, Google Skillshop. Historical credential, expired January 2026.}
\cvitem{Oct 2024}{\textbf{Workover and Well Service Talks}, IATMI Student Chapter ITB. Participant.}
\cvitem{Sep 2024}{\textbf{Automation Clinic: Motor Protection and Control}, Engineering Academy. Participant.}
\cvitem{Sep 2024}{\textbf{How to Write a Good Technical Paper}, SPE. 1.5-hour continuing education session.}
\cvitem{Sep 2024}{\textbf{Geothermal Drilling in Volcanic and Metamorphic Basins}, SPE. 1.5-hour continuing education session.}
\cvitem{Sep 2024}{\textbf{Strategies in Geothermal Exploration: Maximizing Indonesia's Potential}, SPE UNDIP Student Chapter. Guest lecture participant.}
\cvitem{Sep 2024}{\textbf{Ethics of Artificial Intelligence: Navigating Responsible Technology for the Future}, SkilHub Community by Skilvul / Nocturnailed. Webinar participant.}
\cvitem{May 2024}{\textbf{Fundamentals of Digital Marketing}, Google Digital Academy / Skillshop. Completed 27 May 2024.}
\cvitem{May 2024}{\textbf{Introduction to Electronics}, The Open University OpenLearn. Completed a 10-hour course on 21 May 2024.}
\cvitem{May 2024}{\textbf{IT: Device to Device Communication}, The Open University OpenLearn. Completed a 20-hour course on 19 May 2024.}
\cvitem{Mar 2024 - May 2024}{\textbf{College Copilot Bootcamp (CABOOT)}. Completed six sessions, 20 March - 4 May 2024.}
\cvitem{Nov 2020 - Feb 2021}{\textbf{KSN-K Geography Preparation}. Participant in the 8 November 2020 - 7 February 2021 training program.}
% The live profile lists this course without an issue date; do not invent one.
\cvitem{Date not recorded}{\textbf{ML/AI Fundamentals: How Machine Learning Is Shaping the Oil \& Gas Industry}, SPE Java Section.}

\section{Additional Recognition}
\cvitem{May 2024}{\textbf{100 Design Achievement Badge}, Canva.}
\cvitem{2019}{\textbf{Try Out Akbar Winner}, Penerbit Erlangga. Score: 100, SMP/MTs category, Jakarta. SMP Angkasa Halim Perdanakusuma.}

\end{document}
